% ==============================================================================
% ŠABLONA MATURITNÍ PRÁCE S OBHAJOBOU (MPsO)
% Střední průmyslová škola a Gymnázium Na Třebešíně
% ==============================================================================
% Přesná kopie oficiální šablony 2026-27_SPSG_sablona_MPsO.dotx 1:1
% ==============================================================================

\documentclass[12pt, a4paper, oneside]{article}

% ------------------------------------------------------------------------------
% 1. ÚDAJE O PRÁCI (VÝCHOZÍ TEXT ZE ŠABLONY .DOTX)
% ------------------------------------------------------------------------------
\newcommand{\NazevPrace}{Název práce}
\newcommand{\NazevPraceVelkymi}{TÉMA (NÁZEV) MATURITNÍ PRÁCE}
\newcommand{\Autor}{Jméno a příjmení}
\newcommand{\Vedouci}{Jméno Příjmení}
\newcommand{\Trida}{4.}
\newcommand{\SkolniRok}{2025 / 2026}

% ------------------------------------------------------------------------------
% 2. KÓDOVÁNÍ, JAZYK A PÍSMA (Times New Roman 12pt + Arial pro nadpisy)
% ------------------------------------------------------------------------------
\usepackage[utf8]{inputenc}
\usepackage[T1]{fontenc}
\usepackage[czech]{babel}

% Patkové písmo pro běžný text (Times New Roman / TeX Gyre Termes)
\usepackage{tgtermes}
% Bezpatkové písmo pro nadpisy a titulní stranu (Arial / Helvetica / TeX Gyre Heros)
\usepackage{tgheros}
% Neproporcionální písmo pro kód
\usepackage{tgcursor}

% Mikrotypografie
\usepackage{microtype}

% ------------------------------------------------------------------------------
% 3. OKRAJE STRÁNKY A ŘÁDKOVÁNÍ (přesně dle .dotx: okraje 2,5 cm, řádkování 1,5)
% ------------------------------------------------------------------------------
\usepackage[
    a4paper,
    top=2.5cm,
    bottom=2.5cm,
    left=2.5cm,
    right=2.5cm,
    footskip=1.25cm
]{geometry}

% Řádkování 1,5 (odpovídá w:line="360" ve Wordu)
\usepackage{setspace}
\onehalfspacing

% Odstavce bez odsazení prvního řádku, mezera 6pt za odstavcem (w:after="120")
\setlength{\parindent}{0pt}
\setlength{\parskip}{6pt}

% ------------------------------------------------------------------------------
% 4. ZÁHLAVÍ, ZÁPATÍ A ČÍSLOVÁNÍ STRAN (formát "- 5 -" od 5. strany)
% ------------------------------------------------------------------------------
\usepackage{fancyhdr}

% Formát čísla stránky s pomlčkami (odpovídá w:fmt="numberInDash" v šabloně .dotx)
\renewcommand{\thepage}{- \arabic{page} -}

\pagestyle{fancy}
\fancyhf{}
\renewcommand{\headrulewidth}{0pt}
\renewcommand{\footrulewidth}{0pt}
\fancyfoot[C]{\normalfont\fontsize{12pt}{14pt}\selectfont \thepage}

% ------------------------------------------------------------------------------
% 5. FORMÁTOVÁNÍ NADPISŮ (přesně podle stylů Nadpis 1–5 v šabloně .dotx)
% ------------------------------------------------------------------------------
\usepackage{titlesec}
\usepackage{etoolbox}

% Číslování nadpisů až do 4. úrovně (2.1.1.1)
\setcounter{secnumdepth}{4}
% Zobrazení v obsahu do 3. úrovně (2.1.1)
\setcounter{tocdepth}{3}

% Každá hlavní kapitola (\section) začíná na nové stránce
\preto{\section}{\clearpage}

% Nadpis 1 (\section): 16pt, bezpatkové (Arial), tučné, mezera za 12pt
\titleformat{\section}
    {\normalfont\fontsize{16pt}{20pt}\selectfont\sffamily\bfseries}
    {\thesection}{0.75em}{}
\titlespacing*{\section}{0pt}{0pt}{12pt}

% Nadpis 2 (\subsection): 14pt, bezpatkové (Arial), tučné + kurzíva, před 12pt, za 6pt
\titleformat{\subsection}
    {\normalfont\fontsize{14pt}{18pt}\selectfont\sffamily\bfseries\itshape}
    {\thesubsection}{0.75em}{}
\titlespacing*{\subsection}{0pt}{12pt}{6pt}

% Nadpis 3 (\subsubsection): 12pt, bezpatkové (Arial), tučné, před 6pt, za 6pt
\titleformat{\subsubsection}
    {\normalfont\fontsize{12pt}{15pt}\selectfont\sffamily\bfseries}
    {\thesubsubsection}{0.75em}{}
\titlespacing*{\subsubsection}{0pt}{6pt}{6pt}

% Nadpis 4 (\paragraph): 13pt, bezpatkové (Arial), tučné, samostatný řádek
\titleformat{\paragraph}[hang]
    {\normalfont\fontsize{13pt}{16pt}\selectfont\sffamily\bfseries}
    {\theparagraph}{0.75em}{}
\titlespacing*{\paragraph}{0pt}{6pt}{6pt}

% Nadpis 5 / Podnadpis (\subparagraph): 12pt, bezpatkové (Arial), tučné, nečíslovaný
\titleformat{\subparagraph}[hang]
    {\normalfont\fontsize{12pt}{15pt}\selectfont\sffamily\bfseries}
    {}{0pt}{}
\titlespacing*{\subparagraph}{0pt}{6pt}{4pt}

% Nadpis mimo obsah (Čestné prohlášení, Poděkování, Abstrakt, Obsah...)
% V šabloně .dotx: 14pt, bezpatkové (Arial), tučné, nečíslovaný, nezobrazuje se v obsahu
\newcommand{\nadpismimoobsah}[1]{%
    \vspace{6pt}%
    {\normalfont\fontsize{14pt}{18pt}\selectfont\sffamily\bfseries #1\par}%
    \vspace{2pt}%
}

% ------------------------------------------------------------------------------
% 6. OBSAH A SEZNAM OBRÁZKŮ (vodící linky a odsazení dle .dotx)
% ------------------------------------------------------------------------------
\makeatletter
\renewcommand{\@pnumwidth}{3.2em}
\renewcommand{\@tocrmarg}{4.0em}

% Tečkované vodící linky v Obsahu pro \section (stejně jako u podkapitol)
\renewcommand*\l@section[2]{%
  \ifnum \c@tocdepth >\z@
    \addpenalty\@secpenalty
    \addvspace{4pt \@plus\p@}%
    \setlength\@tempdima{1.8em}%
    \begingroup
      \parindent \z@ \rightskip \@pnumwidth
      \parfillskip -\@pnumwidth
      \leavevmode
      \advance\leftskip\@tempdima
      \hskip -\leftskip
      #1\nobreak\leaders\hbox{$\m@th\mkern \@dotsep mu\hbox{.}\mkern \@dotsep mu$}\hfill\nobreak\hb@xt@\@pnumwidth{\hss #2}\par
    \endgroup
  \fi}

% Položky v Seznamu obrázků ("Obrázek 1: Název ...... - 3 -")
\renewcommand*\l@figure[2]{%
  \@dottedtocline{1}{0pt}{5.5em}{Obrázek~#1}{#2}}
\makeatother

% ------------------------------------------------------------------------------
% 7. BALÍČKY PRO GRAFIKU, ODKAZY A POPISKY
% ------------------------------------------------------------------------------
\usepackage{graphicx}
\usepackage{float}
\usepackage{url}

% Popisky obrázků: 12pt, kurzíva, zarovnáno na střed ("Obr. 1: popis")
\usepackage[
    font={it},
    labelfont={it},
    labelsep=colon,
    justification=centering
]{caption}

\addto\captionsczech{%
    \renewcommand{\figurename}{Obr.}%
}

% Klikací odkazy v PDF
\usepackage[hypertexnames=false, hidelinks, unicode]{hyperref}


% ==============================================================================
% DOKUMENT (PŘESNĚ DLE STRAN ŠABLONY .DOTX)
% ==============================================================================
\begin{document}

% Strany 1–4 nemají zobrazené číslo strany
\pagestyle{empty}

% ------------------------------------------------------------------------------
% STRANA 1: TITULNÍ LIST
% ------------------------------------------------------------------------------
\begin{center}
    {\sffamily\fontsize{12pt}{15pt}\selectfont Střední průmyslová škola a~Gymnázium Na Třebešíně\par}
    \vspace{6pt}
    {\sffamily\bfseries\fontsize{16pt}{20pt}\selectfont MATURITNÍ PRÁCE S~OBHAJOBOU\par}

    \vspace{1.2cm}
    \includegraphics[width=5.6cm]{obrazky/logo_skoly.png}

    \vspace{2.2cm}
    {\sffamily\bfseries\fontsize{16pt}{22pt}\selectfont \NazevPraceVelkymi\par}
\end{center}

\vfill

{\sffamily\fontsize{10pt}{14pt}\selectfont
\begin{tabular}{@{}p{3.5cm}l@{}}
    Vypracoval: & \Autor \\[6pt]
    Vedoucí:    & \Vedouci \\[6pt]
    Třída:      & \Trida \\[6pt]
    Školní rok: & \SkolniRok \\
\end{tabular}
}

\vspace{0.5cm}
\clearpage

% ------------------------------------------------------------------------------
% STRANA 2: ZADÁNÍ
% ------------------------------------------------------------------------------
{\sffamily\fontsize{12pt}{15pt}\selectfont VLOŽIT NASKENOVANÉ ZADÁNÍ}

\clearpage

% ------------------------------------------------------------------------------
% STRANA 3: ČESTNÉ PROHLÁŠENÍ
% ------------------------------------------------------------------------------
\vspace*{7.5cm}

\nadpismimoobsah{Čestné prohlášení}

Prohlašuji, že jsem maturitní práci na téma \makebox[3.8cm]{\dotfill} vypracoval(a) samostatně s~využitím odborné literatury a~pramenů uvedených v~přiloženém seznamu využitých zdrojů.

\vspace{1.2cm}
V~Praze dne: 

\vspace{1.8cm}
\textit{Podpis autora}

\clearpage

% ------------------------------------------------------------------------------
% STRANA 4: PODĚKOVÁNÍ
% ------------------------------------------------------------------------------
\vspace*{6.5cm}

\nadpismimoobsah{Poděkování}

Děkuji tímto vedoucímu (vedoucí) maturitní práce panu (paní) \makebox[3.8cm]{\dotfill} za cenné připomínky a~odborné rady při vypracovávání.

\clearpage

% ==============================================================================
% OD STRANY 5: ČÍSLOVÁNÍ STRAN DOLE UPROSTŘED ("- 5 -")
% ==============================================================================
\pagestyle{fancy}
\setcounter{page}{5}

% ------------------------------------------------------------------------------
% STRANA 5: ANOTAČNÍ LIST (ČESKY)
% ------------------------------------------------------------------------------
\nadpismimoobsah{Název práce:}
Název práce

\nadpismimoobsah{Autor:}
Autor práce

\nadpismimoobsah{Vedoucí práce:}
Vedoucí práce

\nadpismimoobsah{Abstrakt}
Krátký souhrn, o čem práce pojednává

\nadpismimoobsah{Klíčová slova}
Klíčová slova

\clearpage

% ------------------------------------------------------------------------------
% STRANA 6: ANOTAČNÍ LIST (ANGLICKY)
% ------------------------------------------------------------------------------
\nadpismimoobsah{Summary}
Souhrn přeložený do AJ 

\nadpismimoobsah{Keywords}
Klíčová slova v angličtině

\clearpage

% ------------------------------------------------------------------------------
% STRANA 7: OBSAH
% ------------------------------------------------------------------------------
\nadpismimoobsah{Obsah}
\makeatletter
\@starttoc{toc}
\makeatother

\clearpage

% ------------------------------------------------------------------------------
% STRANA 8+: TEXT PRÁCE (PŘESNĚ DLE ŠABLONY .DOTX)
% ------------------------------------------------------------------------------
\section{Úvod}
Text úvodu…

\begin{figure}[H]
    \centering
    % Placeholder pro obrázek
    \caption{např. Úvodní schéma}
\end{figure}

\section{Kapitola}

\subsection{Kapitola}

\subsubsection{Kapitola}

\paragraph{Kapitola}

\subparagraph{Podnadpis}

\clearpage

\section{Závěr}
Závěr…

\clearpage

\section{Seznam obrázků}
\makeatletter
\@starttoc{lof}
\makeatother

\clearpage

\section{Seznam využitých zdrojů}

\end{document}
